\centerline{\textbf{\large Price Dynamics in the Construction Materials}}
\centerline{\textbf{\large Sector: Patterns and Data-Driven Predictions }}
\vspace{2cm}

\rightline{Author: Diego Milciades Espínola Alvarenga}
\rightline{Matías Rubén Sánchez Sánchez}
\vspace{0.5cm}
\rightline{Advisor: D.Sc. Gustavo Sosa-Cabrera}

\vspace{0.5cm}

\noindent
\textbf{ABSTRACT}
\justify
Construction material prices in Paraguay fluctuate in ways that are difficult to anticipate, yet these fluctuations directly affect project budgets, procurement planning, and overall economic viability. This study addresses that challenge by developing predictive models for three widely consumed materials — Yguazú cement, Tobatí common brick, and Inpacal 20 kg hydrated lime — using monthly price records from January 2014 to July 2025. Two exogenous variables were incorporated: the COVID-19 lockdown period and the daily water level of the Paraguay River at the Port of Asunción, whose 2021 historic low forced the National Cement Industry to shift from river to road transport, raising distribution costs substantially.

\justify
Missing price values were handled through quadratic polynomial interpolation, chosen after comparing several methods against verified reference data from a local supplier. From there, three forecasting approaches were built: a SARIMAX statistical model and two recurrent neural network architectures, LSTM and GRU. All models were tested over a 24-month horizon under alternative exogenous-variable scenarios, with RMSE as the primary evaluation metric.

\justify
Neural network models outperformed the statistical baseline. The multivariate LSTM achieved the highest overall accuracy, while GRU reached comparable results at lower computational cost. SARIMAX proved useful as a reference point but struggled with the nonlinear dynamics present in the series. A key finding was that forecasting river levels in advance with a univariate LSTM — and then feeding those projections into the price models as an exogenous input — measurably improved prediction accuracy, supporting the hypothesis that hydrological conditions influence construction material costs in Paraguay.

\justify
The models were built entirely with publicly available data and open-source tools. Beyond the immediate application, this work establishes a reproducible framework that can be extended to other materials or economic indicators relevant to the Paraguayan construction sector.

\vspace{0.5cm}

\noindent
\textbf{Keywords}: Construction material prices, exogenous variables, Paraguay River, COVID-19, predictive models, SARIMAX, LSTM, GRU, Paraguay.

